\section{Background}
\label{sec:background}

\subsection{\vpcc and the Delivery Unit}
\label{sec:bg-codecs}

A dynamic point cloud is a sequence of unordered 3D points with colour
attributes, viewed with \gls{6dof}; the \gls{mpeg}/8i human sequences
used here carry about $10^6$ points per frame at 30\,fps and require
gigabits per second uncompressed~\cite{8i,pcstreaming21}. \vpcc, the
principal component of the \gls{v3c} standard ISO/IEC
23090-5~\cite{v3c,vpccgraziosi}, compresses such content by projection:
the cloud is segmented into connected patches, each patch is projected
onto a 2D plane, and the projections are packed into atlases coded as
three ordinary video streams (geometry, attribute and a lossless
occupancy map) with \gls{hevc}~\cite{hevc} in the reference software
\gls{tmc2}~\cite{tmc2}. The decoder back-projects the decoded pixels
into 3D.

Because the maps are ordinary video, \vpcc inherits the video codec's
prediction structure, and this is what matters for transport. Frames
are coded in \glspl{gof}. The \gls{ra} configuration of the \vpcc common
test conditions~\cite{vpccctc} uses I-, P- and B-frames with a \gof of
32, maximising compression at the cost of long dependency chains; the
\gls{ldp} configuration uses only I- and P-frames with \gofs of 16 or 4,
trading compression for resilience~\cite{hmctc}. In either case a \gof
must be received in full before any of its frames can be reconstructed,
and in a streaming deployment each \gof carries a presentation
deadline. The \gof is therefore the \emph{delivery unit} of a \vpcc
stream, and its size is the structural variable of \crystal: a smaller
\gof puts fewer packets at risk per unit but leaves a tighter deadline
budget.

\subsection{Reliable \quic Delivery}
\label{sec:bg-quic}

Volumetric video, like 2D video at scale, is delivered as time-bounded
HTTP objects: ISO/IEC 23090-10 packages \vpcc into ISOBMFF tracks and
\gls{dash} segments~\cite{v3ccarriage}, and the 5G-MAG platform plays
them from an HTTP origin~\cite{fivegmag}. Our testbed keeps this object
model and the 5G-MAG decoder but replaces the \gls{dash} client with a
native \gls{http3} fetcher, so that each \gof is one HTTP object carried
over \quic~\cite{mosaic2026}. Over \gls{tcp}, a lost segment blocks every
byte behind it on the connection, so a loss inside one \gof also stalls
the next \gof in flight. \quic~\cite{rfc9000}, the substrate of
\gls{http3}~\cite{rfc9114}, removes that coupling with independent
streams, combines the \gls{tls}~1.3 and transport handshakes into one
round trip, and is deployed widely: in October 2026, 40.9\% of public
web sites advertised \gls{http3}~\cite{w3techs2026}.

What matters for \crystal is how standard \quic recovers a loss. Packet
numbers are never reused, so the sender detects a loss either when a
packet at least $k_{\mathrm{PacketThreshold}}=3$ later is acknowledged
or, when no later acknowledgement arrives, when a \gls{pto} expires~\cite{rfc9002}:
\begin{equation}
  \mathrm{PTO} = \mathrm{sRTT} + \max(4\,\mathrm{RTT}_{\mathrm{var}},
  k_{\mathrm{granularity}}) + \mathrm{max\_ack\_delay}.
  \label{eq:pto}
\end{equation}
A loss therefore costs at least one round trip, and a full \gls{pto}
when the lost packet is among the last of a \gof (a tail loss), and
each recovery is followed by a congestion-window reduction that
persists into later transfers. Because \quic is reliable, the decoder
never receives an incomplete \gof: loss becomes \emph{delay}, and delay
becomes a \emph{stall} when the recovered \gof misses its deadline.
Measurements on our testbed show the two impairments leave distinct
signatures, added latency shifting the whole delivery-time distribution
and loss expanding its tail~\cite{vquest2026}, and the per-\gof
deadline-miss rate correlates only partially with subjective quality
($\rho = -0.34$) because a missed \gof of 4 frames and one of 16 frames
remove different durations of content~\cite{mosaic2026}. The temporal
loss fraction of Section~\ref{sec:model} is defined to close that gap.
